\chapter{Trigonometria: o círculo trigonométrico}\label{ch:g11:trigo}

A trigonometria começa nos triângulos retângulos, mas sua verdadeira casa
é o \emph{círculo trigonométrico}, onde o \omterm{def:g11:trigo:cossin}{cosseno} e o \omterm{def:g11:trigo:cossin}{seno} passam a ser
funções de um \omterm{def:g10:numbers:sets}{número real} qualquer. Este capítulo instala o \omterm{def:g11:trigo:radian}{radiano}, o
enrolamento da reta real em torno do círculo, os valores exatos que se
devem saber de cor e as simetrias que geram todas as identidades
clássicas. O estudo de $\cos$ e $\sin$ como funções --- variação,
\omterm{def:g11:deriv:derivative}{derivadas} --- é feito no \cref{ch:g12:trigo}.

\section{Radianos e o enrolamento}

\begin{definition}[Radiano]\label{def:g11:trigo:radian}
Seja $\mathcal C$ o círculo de raio $1$ centrado na origem $O$, e seja
$I(1, 0)$. A medida em \emph{radianos}\index{radiano} de um ângulo
central de $\mathcal C$ é o comprimento do arco que ele recorta sobre
$\mathcal C$. Como o círculo completo tem comprimento $2\pi$:
\[
360^\circ = 2\pi \text{ rad},
\qquad
180^\circ = \pi \text{ rad},
\qquad
1 \text{ rad} = \frac{180^\circ}{\pi} \approx 57.3^\circ .
\]
\end{definition}

\begin{method}[Converter]\label{met:g11:trigo:convert}
Graus e \omterm{def:g11:trigo:radian}{radianos} são proporcionais: multiplique por $\frac{\pi}{180}$
para ir de graus a \omterm{def:g11:trigo:radian}{radianos} e por $\frac{180}{\pi}$ no sentido inverso.
Por exemplo, $60^\circ = 60 \times \frac{\pi}{180} = \frac{\pi}{3}$ e
$\frac{3\pi}{4} = \frac{3 \times 180^\circ}{4} = 135^\circ$.
\end{method}

\begin{definition}[Enrolar a reta no círculo]\label{def:g11:trigo:winding}
A cada \omterm{def:g10:numbers:sets}{número real} $t$ associe o ponto $M(t)$ de $\mathcal C$ obtido
percorrendo uma distância $\abs t$ ao longo do círculo a partir de $I$,
no sentido anti-horário se $t \geq 0$ e no horário se $t < 0$. Como o
comprimento do círculo é $2\pi$, os números $t$ e $t + 2k\pi$
($k \in \Z$) chegam ao mesmo ponto: $M(t + 2k\pi) = M(t)$.
\end{definition}

\begin{example}
$M(0) = I$; $M\!\left(\frac{\pi}{2}\right) = (0, 1)$, o topo do círculo;
$M(\pi) = (-1, 0)$; $M\!\left(-\frac{\pi}{2}\right) = (0, -1)$; e
$M\!\left(\frac{9\pi}{4}\right) = M\!\left(\frac{\pi}{4} +
2\pi\right) = M\!\left(\frac{\pi}{4}\right)$.
\end{example}

\section{Cosseno e seno}

\begin{definition}[Cosseno e seno]\label{def:g11:trigo:cossin}
Para todo real $t$, o \emph{cosseno}\index{cosseno} e o
\emph{seno}\index{seno} de $t$ são as \omterm{def:g10:coordgeom:system}{coordenadas} do ponto $M(t)$ do
círculo trigonométrico:
\[
M(t) = (\cos t,\ \sin t).
\]
\end{definition}

\begin{omfigure}
\begin{tikzpicture}[scale=2.0]
  \draw[->] (-1.3,0) -- (1.4,0) node[right] {$x$};
  \draw[->] (0,-1.3) -- (0,1.4) node[above] {$y$};
  \draw[gray] (0,0) circle (1);
  \coordinate (O) at (0,0);
  \coordinate (M) at (55:1);
  \draw[thick,omDef] (O) -- (M);
  \draw[dashed] (M) -- (M |- O) node[below, font=\small] {$\cos t$};
  \draw[dashed] (M) -- (M -| O) node[left, font=\small] {$\sin t$};
  \draw[->,omThm,thick] (0.35,0) arc (0:55:0.35);
  \node[omThm, font=\small] at (27:0.55) {$t$};
  \fill (M) circle (0.7pt) node[above right] {$M(t)$};
  \fill (1,0) circle (0.7pt) node[below right] {$I$};
\end{tikzpicture}

{\small O círculo trigonométrico: o \omterm{def:g10:numbers:sets}{número real} $t$, enrolado no sentido
anti-horário a partir de $I$, chega ao ponto $M(t)$, cujas \omterm{def:g10:coordgeom:system}{coordenadas}
\emph{definem} $\cos t$ e $\sin t$.}
\end{omfigure}

\begin{proposition}[Primeiras propriedades]\label{prop:g11:trigo:first}
Para todos $t \in \R$ e $k \in \Z$:
\[
-1 \leq \cos t \leq 1, \qquad
-1 \leq \sin t \leq 1, \qquad
\cos^2 t + \sin^2 t = 1,
\]
\[
\cos(t + 2k\pi) = \cos t, \qquad \sin(t + 2k\pi) = \sin t .
\]
\end{proposition}

\begin{proof}
$M(t)$ está no círculo de raio $1$, de modo que suas \omterm{def:g10:coordgeom:system}{coordenadas} estão
entre $-1$ e $1$ e satisfazem a \omterm{def:g10:algebra:equation}{equação} do círculo $x^2 + y^2 = 1$, que é
a identidade $\cos^2 t + \sin^2 t = 1$. A periodicidade reenuncia
$M(t + 2k\pi) = M(t)$ (\cref{def:g11:trigo:winding}).
\end{proof}

\begin{example}\label{ex:g11:trigo:identity}
Se $\sin t = \frac35$ e $t \in \intcc{\frac\pi2}{\pi}$ (segundo
quadrante), então $\cos^2 t = 1 - \frac{9}{25} = \frac{16}{25}$, logo
$\cos t = \pm\frac45$; no segundo quadrante a \omterm{def:g10:coordgeom:system}{abscissa} é negativa,
portanto $\cos t = -\frac45$.
\end{example}

Os valores que se devem saber de cor:
\begin{center}
\begin{tabular}{c|ccccc}
$t$ & $0$ & $\dfrac{\pi}{6}$ & $\dfrac{\pi}{4}$ & $\dfrac{\pi}{3}$ & $\dfrac{\pi}{2}$\\[6pt]
\hline\rule{0pt}{14pt}%
$\cos t$ & $1$ & $\dfrac{\sqrt3}{2}$ & $\dfrac{\sqrt2}{2}$ & $\dfrac12$ & $0$\\[6pt]
$\sin t$ & $0$ & $\dfrac12$ & $\dfrac{\sqrt2}{2}$ & $\dfrac{\sqrt3}{2}$ & $1$
\end{tabular}
\end{center}

\begin{remark}
Um recurso de memória: os \omterm{def:g11:trigo:cossin}{senos} se leem
$\frac{\sqrt0}{2}, \frac{\sqrt1}{2}, \frac{\sqrt2}{2}, \frac{\sqrt3}{2},
\frac{\sqrt4}{2}$, e os \omterm{def:g11:trigo:cossin}{cossenos} são a mesma lista ao contrário.
\end{remark}

\section{Ângulos associados}

\begin{proposition}[Ângulos associados]\label{prop:g11:trigo:associated}
\index{ângulos associados}
Para todo $t \in \R$:
\[
\begin{aligned}
\cos(-t) &= \cos t, & \sin(-t) &= -\sin t,\\
\cos(\pi - t) &= -\cos t, & \sin(\pi - t) &= \sin t,\\
\cos(\pi + t) &= -\cos t, & \sin(\pi + t) &= -\sin t,\\
\cos\left(\tfrac{\pi}{2} - t\right) &= \sin t, &
\sin\left(\tfrac{\pi}{2} - t\right) &= \cos t .
\end{aligned}
\]
\end{proposition}

\begin{proof}
Cada linha exprime uma simetria do círculo.

$M(-t)$ é o reflexo de $M(t)$ no eixo $x$: a \omterm{def:g10:coordgeom:system}{abscissa} não muda, a
ordenada troca de sinal.

$M(\pi - t)$ é o reflexo de $M(t)$ no eixo $y$: andar para trás a partir
de $\pi$ por $t$ leva ao ponto oposto (na horizontal) ao de andar para a
frente a partir de $0$ por $t$. A \omterm{def:g10:coordgeom:system}{abscissa} troca de sinal, a ordenada não
muda.

$M(\pi + t)$ é diametralmente oposto a $M(t)$ (meia volta): as duas
\omterm{def:g10:coordgeom:system}{coordenadas} trocam de sinal.

$M\!\left(\frac\pi2 - t\right)$ é o reflexo de $M(t)$ na reta diagonal
$y = x$, que troca as duas \omterm{def:g10:coordgeom:system}{coordenadas}.
\end{proof}

\begin{omfigure}
\begin{tikzpicture}[scale=2.0]
  \draw[->] (-1.35,0) -- (1.45,0) node[right] {$x$};
  \draw[->] (0,-1.35) -- (0,1.45) node[above] {$y$};
  \draw[gray] (0,0) circle (1);
  \fill (40:1) circle (0.7pt) node[above right, font=\small] {$M(t)$};
  \fill (-40:1) circle (0.7pt) node[below right, font=\small] {$M(-t)$};
  \fill (140:1) circle (0.7pt) node[above left, font=\small]
    {$M(\pi - t)$};
  \fill (220:1) circle (0.7pt) node[below left, font=\small]
    {$M(\pi + t)$};
  \draw[black, dashed, thin] (40:1) -- (-40:1);
  \draw[black, dashed, thin] (40:1) -- (140:1);
  \draw[black, dashed, thin] (140:1) -- (220:1);
  \draw[omDef, thin] (0,0) -- (40:1);
  \draw[omDef, thin] (0,0) -- (140:1);
  \draw[omDef, thin] (0,0) -- (220:1);
  \draw[omDef, thin] (0,0) -- (-40:1);
\end{tikzpicture}

{\small Os quatro pontos associados: $M(-t)$ espelha $M(t)$ no eixo $x$,
$M(\pi - t)$ no eixo $y$, e $M(\pi + t)$ é diametralmente oposto. Toda
identidade da \cref{prop:g11:trigo:associated} se lê nesta figura.}
\end{omfigure}

\begin{example}
$\cos\frac{2\pi}{3} = \cos\left(\pi - \frac\pi3\right)
= -\cos\frac\pi3 = -\frac12$, e
$\sin\frac{7\pi}{6} = \sin\left(\pi + \frac\pi6\right)
= -\sin\frac\pi6 = -\frac12$. A tabela de cinco valores, estendida pelas
simetrias, cobre o círculo inteiro.
\end{example}

\section{Resolver equações trigonométricas}

\begin{theorem}[As equações $\cos t = \cos a$ e $\sin t = \sin a$]
\label{thm:g11:trigo:equations}
Para \omterm{def:g10:numbers:sets}{números reais} $t$ e $a$:
\[
\cos t = \cos a \iff t = a + 2k\pi \ \text{ou}\ t = -a + 2k\pi
\quad (k \in \Z),
\]
\[
\sin t = \sin a \iff t = a + 2k\pi \ \text{ou}\ t = \pi - a + 2k\pi
\quad (k \in \Z).
\]
\end{theorem}

\begin{proof}
Dois pontos do círculo trigonométrico têm a mesma \omterm{def:g10:coordgeom:system}{abscissa} exatamente
quando são iguais ou reflexos um do outro no eixo $x$; pela
\cref{prop:g11:trigo:associated}, o reflexo de $M(a)$ é $M(-a)$. Assim,
$\cos t = \cos a$ significa $M(t) = M(a)$ ou $M(t) = M(-a)$, isto é,
$t = \pm a$ a menos de um múltiplo de $2\pi$. Analogamente, ordenadas
iguais significam pontos iguais ou reflexos no eixo $y$, e o reflexo de
$M(a)$ é $M(\pi - a)$.
\end{proof}

\begin{method}[Resolver $\cos t = c$ em um intervalo]\label{met:g11:trigo:solve}
\begin{enumerate}
  \item Encontre \emph{um} ângulo $a$ com $\cos a = c$, a partir da
        tabela de valores conhecidos.
  \item Escreva as soluções gerais $t = \pm a + 2k\pi$
        (\cref{thm:g11:trigo:equations}).
  \item Escolha os inteiros $k$ que caem dentro do \omterm{def:g10:numbers:interval}{intervalo} pedido e
        liste as soluções.
\end{enumerate}
Proceda do mesmo modo para $\sin t = c$, com $t = a + 2k\pi$ ou
$t = \pi - a + 2k\pi$.
\end{method}

\begin{example}\label{ex:g11:trigo:solve}
Resolva $\cos t = \frac12$ em $\intoc{-\pi}{\pi}$. Uma solução de
$\cos a = \frac12$ é $a = \frac{\pi}{3}$. As soluções gerais são
$t = \frac\pi3 + 2k\pi$ e $t = -\frac\pi3 + 2k\pi$. Dentro de
$\intoc{-\pi}{\pi}$, só $k = 0$ contribui:
$t \in \left\{-\frac\pi3,\ \frac\pi3\right\}$.

Resolva $\sin t = -\frac{\sqrt2}{2}$ em $\intco{0}{2\pi}$. Um ângulo é
$a = -\frac\pi4$; as soluções são $t = -\frac\pi4 + 2k\pi$ e
$t = \pi + \frac\pi4 + 2k\pi = \frac{5\pi}{4} + 2k\pi$. Em
$\intco{0}{2\pi}$: $t = \frac{7\pi}{4}$ (de $k = 1$) e
$t = \frac{5\pi}{4}$.
\end{example}

\begin{omfigure}
\begin{tikzpicture}[scale=2.0]
  \draw[->] (-1.35,0) -- (1.45,0) node[right] {$x$};
  \draw[->] (0,-1.2) -- (0,1.3) node[above] {$y$};
  \draw[gray] (0,0) circle (1);
  \draw[omProp, thick] (0.5,-1.05) -- (0.5,1.15)
    node[above, font=\small] {$x = \tfrac12$};
  \fill (60:1) circle (0.7pt)
    node[above right, font=\small] {$M\!\left(\tfrac\pi3\right)$};
  \fill (-60:1) circle (0.7pt)
    node[below right, font=\small] {$M\!\left(-\tfrac\pi3\right)$};
  \draw[omDef, thin] (0,0) -- (60:1);
  \draw[omDef, thin] (0,0) -- (-60:1);
\end{tikzpicture}

{\small Resolver $\cos t = \frac12$: a reta vertical $x = \frac12$
(laranja) corta o círculo trigonométrico em dois pontos, simétricos em
relação ao eixo $x$ --- daí as duas famílias de soluções
$t = \pm\frac\pi3 + 2k\pi$.}
\end{omfigure}

\section{Exercícios}

\begin{exercise}[$\star$]\label{exo:g11:trigo:1}
Converta para \omterm{def:g11:trigo:radian}{radianos}: $30^\circ$, $45^\circ$, $120^\circ$,
$270^\circ$. Converta para graus: $\frac{\pi}{5}$, $\frac{5\pi}{6}$,
$\frac{7\pi}{4}$, $\frac{2\pi}{9}$.
\end{exercise}

\begin{exercise}[$\star$]\label{exo:g11:trigo:2}
Marque no círculo trigonométrico os pontos $M(t)$ para
\[
t = \frac{2\pi}{3},\quad
t = -\frac{\pi}{4},\quad
t = \frac{17\pi}{6},\quad
t = -\frac{7\pi}{2},
\]
depois de reduzir cada um a um valor de $\intoc{-\pi}{\pi}$ módulo
$2\pi$.
\end{exercise}

\begin{exercise}[$\star$]\label{exo:g11:trigo:3}
Usando a tabela e os \omterm{prop:g11:trigo:associated}{ângulos associados}, dê os valores exatos de
\[
\cos\frac{3\pi}{4}, \quad
\sin\frac{5\pi}{6}, \quad
\cos\left(-\frac{\pi}{3}\right), \quad
\sin\frac{4\pi}{3}, \quad
\cos\frac{11\pi}{6} .
\]
\end{exercise}

\begin{exercise}[$\star$]\label{exo:g11:trigo:4}
Dado $\cos t = \frac{5}{13}$ com $t \in \intoo{-\frac\pi2}{0}$, calcule
$\sin t$ exatamente.
\end{exercise}

\begin{exercise}[$\star\star$]\label{exo:g11:trigo:5}
Simplifique, para todo real $x$:
\[
A = \cos(\pi + x) + \cos(-x) + \sin\left(\frac\pi2 - x\right) - \cos x,
\qquad
B = \sin(\pi - x) + \sin(\pi + x) + \sin(-x) .
\]
\end{exercise}

\begin{exercise}[$\star\star$]\label{exo:g11:trigo:6}
Resolva em $\intoc{-\pi}{\pi}$:
\[
\cos t = \frac{\sqrt3}{2}, \qquad
\sin t = \frac12, \qquad
\cos t = -1 .
\]
\end{exercise}

\begin{exercise}[$\star\star$]\label{exo:g11:trigo:7}
Resolva em $\intco{0}{2\pi}$:
\[
\sin t = -\frac{\sqrt3}{2}, \qquad
\cos t = 0, \qquad
2\sin t + 1 = 0 .
\]
\end{exercise}

\begin{exercise}[$\star\star$]\label{exo:g11:trigo:8}
Resolva em $\R$ a \omterm{def:g10:algebra:equation}{equação} $\cos 2t = \cos t$. (Use o
\cref{thm:g11:trigo:equations} com $a = t$ e discuta as duas famílias de
soluções.)
\end{exercise}

\begin{exercise}[$\star\star$]\label{exo:g11:trigo:9}
Mostre que, para todo real $x$,
\[
\bigl(\cos x + \sin x\bigr)^2 + \bigl(\cos x - \sin x\bigr)^2 = 2 .
\]
\end{exercise}

\begin{exercise}[$\star\star$]\label{exo:g11:trigo:10}
Uma roda de raio $30$ cm rola sem deslizar. De que ângulo (em \omterm{def:g11:trigo:radian}{radianos} e
depois em graus) ela gira quando a bicicleta avança $1.5$ m? Quanto a
bicicleta avança durante uma volta completa da roda?
\end{exercise}

\begin{exercise}[$\star\star\star$]\label{exo:g11:trigo:11}
Resolva em $\intoc{-\pi}{\pi}$ a inequação $\cos t \leq \frac12$.
(Esboce o círculo trigonométrico, marque a região em que a \omterm{def:g10:coordgeom:system}{abscissa} é no
máximo $\frac12$ e leia o \omterm{def:g10:numbers:interval}{intervalo} de valores de $t$.)
\end{exercise}

\section{Problema: A roda-gigante e a maré}

\begin{problem}[{Problema de fim de semana --- o radiano como
distância rolada e o círculo trigonométrico como relógio-mestre de
todo fenômeno periódico}]
\label{pb:g11:trigo:1}
Uma roda-gigante leva você em torno de um círculo; a maré leva a
água do porto para cima e para baixo pelo mesmo círculo
matemático. Todo fenômeno periódico --- rodas, marés, som,
estações --- lê seu horário no círculo trigonométrico deste
capítulo (\cref{def:g11:trigo:cossin}), e as \omterm{def:g10:algebra:equation}{equações} do
\cref{met:g11:trigo:solve} dizem os instantes exatos. Este
problema converte, enrola, modela e resolve --- e decide de
passagem por que os matemáticos medem ângulos em \omterm{def:g11:trigo:radian}{radianos}.

\medskip
\noindent\textbf{Parte I --- \omterm{def:g11:trigo:radian}{Radianos}, a distância rolada.}
\begin{enumerate}
  \item Converta para \omterm{def:g11:trigo:radian}{radianos}: $30^\circ$, $135^\circ$,
        $300^\circ$; e para graus: $\frac{3\pi}{4}$,
        $\frac{5\pi}{6}$ (\cref{met:g11:trigo:convert}).
  \item A razão de ser dos \omterm{def:g11:trigo:radian}{radianos}: um arco de ângulo $t$ (em
        \omterm{def:g11:trigo:radian}{radianos}) num círculo de raio $r$ tem comprimento
        exatamente $r\,t$. Que arco um ângulo de $2$ \omterm{def:g11:trigo:radian}{radianos}
        recorta num círculo de raio $3$ m? (Nenhum $\pi$ em lugar
        algum --- é exatamente esse o ponto.)
  \item Uma roda de raio $30$ cm rola $1$ km sem deslizar. De
        quantos \emph{\omterm{def:g11:trigo:radian}{radianos}} ela girou, e quantas voltas
        completas isso representa?
  \item Marque no círculo trigonométrico e calcule exatamente:
        $\cos\frac{2\pi}{3}$, $\sin\left(-\frac{\pi}{4}\right)$ e
        $\cos\frac{19\pi}{6}$ (reduza módulo $2\pi$ primeiro;
        \cref{prop:g11:trigo:associated}).
  \item Um ângulo $t \in \intoo{\frac\pi2}{\pi}$ satisfaz
        $\sin t = \frac35$. Usando $\cos^2 t + \sin^2 t = 1$ e o
        quadrante, encontre $\cos t$ e $\tan t$.
\end{enumerate}

\medskip
\noindent\textbf{Parte II --- A roda-gigante.}
Uma roda-gigante tem raio $60$ m, o eixo a $65$ m do chão, e dá
uma volta a cada $30$ minutos. Você embarca no ponto mais baixo no
instante $t = 0$ (em minutos), de modo que sua altura obedece a
\[
h(t) = 65 - 60 \cos\!\left(\frac{\pi t}{15}\right).
\]
\begin{enumerate}[resume]
  \item Justifique a fórmula: confira o ângulo girado após $t$
        minutos, a altura em $t = 0$ e explique o sinal de menos.
  \item Calcule sua altura em $t = 7.5$, $t = 15$ e $t = 20$
        minutos (valores exatos).
  \item Em que instante você chega \emph{pela primeira vez} a
        $95$ m? (Resolva $h(t) = 95$ com o
        \cref{met:g11:trigo:solve}.)
  \item Em que \omterm{def:g10:numbers:interval}{intervalo} de tempo do passeio você fica a pelo
        menos $95$ m de altura --- e por quanto tempo ao todo?
        (Compare com o \cref{exo:g11:trigo:11}.)
  \item Onde a altura muda \emph{mais depressa}? Compare a
        subida durante o primeiro minuto ($h(1) - h(0)$) com a
        subida entre os minutos $7$ e $8$, e formule a
        observação (a ferramenta que a torna exata --- a \omterm{def:g11:deriv:derivative}{derivada}
        da senoide --- chega no ano 12).
  \item Vez do engenheiro: projete a fórmula da altura de uma
        roda com altura máxima $40$ m, altura mínima $4$ m e
        período $12$ minutos, embarcando embaixo em $t = 0$.
\end{enumerate}

\medskip
\noindent\textbf{Parte III --- A maré e as simetrias do
círculo.}
Em um porto, a profundidade da água em metros, $t$ horas após a
meia-noite, é
\[
d(t) = 7 + 3 \sin\!\left(\frac{\pi t}{6}\right).
\]
\begin{enumerate}[resume]
  \item Dê as profundidades máxima e mínima, os instantes em que
        elas ocorrem pela primeira vez e o período da maré.
  \item Um navio cargueiro precisa de $8.5$ m de água. Resolva
        $d(t) \geq 8.5$ em $\intco{0}{12}$: em que janela ele
        pode entrar e qual é a duração dessa janela? Quando a
        janela seguinte se abre?
  \item Rededuza duas identidades de \omterm{prop:g11:trigo:associated}{ângulos associados}
        diretamente das simetrias do círculo ---
        $\sin(\pi - t) = \sin t$ (espelho no eixo vertical) e
        $\cos(\pi + t) = -\cos t$ (meia volta) --- e enuncie a
        identidade complementar
        $\cos\left(\frac\pi2 - t\right) = \sin t$.
  \item Paridade (vocabulário do \cref{pb:g11:func:1}):
        classifique $\sin$, $\cos$ e $t \mapsto \sin t + t^3$
        como pares ou ímpares, com demonstrações de uma linha a
        partir do círculo.
  \item Resolva em $\intco{0}{2\pi}$: $2\sin^2 t - \sin t - 1 =
        0$. (Um \omterm{def:g11:quad:quadratic}{trinômio do segundo grau} disfarçado: fatore-o
        como no \cref{pb:g11:quad:1} e depois leia o círculo.)
\end{enumerate}

\medskip
\noindent\textbf{Parte IV --- O ângulo do matemático.}
\begin{enumerate}[resume]
  \item Quantos graus tem $1$ \omterm{def:g11:trigo:radian}{radiano}? E uma armadilha da
        calculadora: qual é maior, $\sin(1)$ (modo \omterm{def:g11:trigo:radian}{radiano}) ou
        $\sin(1^\circ)$ --- e por que fator, mais ou menos? Que
        moral fica para a configuração da calculadora?
  \item Resolva $\cos t = \sin t$ em $\intco{0}{2\pi}$, usando o
        círculo (onde \omterm{def:g10:coordgeom:system}{abscissa} e ordenada são iguais?).
  \item Para ângulos pequenos, o arco e seu \omterm{def:g11:trigo:cossin}{seno} quase
        coincidem: compare $\sin(0.1)$ com $0.1$ (cinco casas
        decimais) e dê o erro relativo. Essa ``\omterm{def:g10:numbers:approx}{aproximação} de
        ângulos pequenos'' governa navios e pêndulos; o teorema
        por trás dela ($\frac{\sin t}{t} \to 1$) é uma estrela do
        capítulo de limites do ano 12.
  \item Final --- o círculo trigonométrico como relógio-mestre,
        uma frase por item: os \omterm{def:g11:trigo:radian}{radianos} transformam ângulos em
        números honestos (arco $=$ raio $\times$ ângulo); \omterm{def:g11:trigo:cossin}{cosseno}
        e \omterm{def:g11:trigo:cossin}{seno} são a sombra e a altura do ponteiro do relógio; as
        simetrias do círculo são as identidades; resolver
        $\cos t = c$ lê horários no relógio; e todo fenômeno
        periódico --- roda, maré, som --- é esse relógio
        fantasiado.
\end{enumerate}
\end{problem}
